\chapter{Multi-armed bandits}
\label{ch:bandits}

Before facing sequential decisions, we study the exploration problem in its
purest form. A \emph{multi-armed bandit} is an MDP with a single state:
each action (``arm'') gives a random reward drawn from an unknown
distribution, and the agent must find the best arm while collecting as much
reward as possible along the way. Every reward spent trying an arm is a
reward not collected from the best one; every arm not tried may be the
best. This \emph{exploration--exploitation dilemma} reappears in all of RL
(the $\varepsilon$-greedy policy of DQN, the entropy bonus of PPO,
\cref{ch:exploration}), and the bandit algorithms of this chapter are used
directly in planning: Monte Carlo tree search chooses its simulations with
the UCB rule (\cref{ch:model-based}).

\section{The problem}

\begin{definition}{Stochastic multi-armed bandit}{bandit}
  A bandit with $K$ arms is a family of reward distributions $\nu_1, \dots,
  \nu_K$ on $\R$ with means $\mu_1, \dots, \mu_K$. At each round $t = 1, 2,
  \dots, T$ the agent chooses an arm $A_t$ and receives a reward $R_t \sim
  \nu_{A_t}$, independent of the past given $A_t$. Let $\mu^\star = \max_a
  \mu_a$, $\Delta_a = \mu^\star - \mu_a$ the \emph{gap} of arm $a$, and
  $N_a(t)$ the number of times arm $a$ was played in the first $t$ rounds.
  The \emph{regret} after $T$ rounds is
  \begin{equation}
    \mathcal{R}_T = T \mu^\star - \E\left[ \sum_{t=1}^{T} R_t \right].
  \end{equation}
\end{definition}

\begin{intuition}
  $T \mu^\star$ is what an oracle who knew the best arm from the start
  would collect on average in $T$ rounds; $\E[\sum_t R_t]$ is what the
  agent collects. The regret is the price of not knowing which arm is best.
  It compares with the best arm's \emph{mean}, not with the luckiest
  rewards, so it measures bad decisions, not bad luck. If $\mathcal{R}_T$
  grows linearly in $T$, the agent keeps losing a fixed amount per round:
  it never learns. If it grows sublinearly, the loss per round
  $\mathcal{R}_T / T$ tends to $0$: the agent ends up playing the best arm.
\end{intuition}

\begin{remark}{The regret as an expected loss}{regret-loss}
  Imagine that, before the game, nature draws a whole table of rewards:
  for every round $t$ and every arm $a$, a reward $X_{a,t} \sim \nu_a$, all
  independent. At round $t$ the agent receives the entry of the arm it
  plays, $R_t = X_{A_t, t}$, and never sees the others. (This gives the
  same game as \cref{def:bandit}.) Let $\star$ be a best arm. The loss of
  round $t$ is $X_{\star, t} - X_{A_t, t}$, what the best arm would have
  given minus what the agent got, and since $\E[X_{\star, t}] = \mu^\star$,
  \begin{equation}
    \mathcal{R}_T = \E\left[ \sum_{t=1}^{T} \left( X_{\star, t} - X_{A_t, t}
    \right) \right]:
  \end{equation}
  the regret is the expected total loss with respect to the best arm.

  The loss of one round can be negative: a bad arm sometimes gets lucky.
  Its \emph{expected} value is not, for the following reason. The choice
  of the agent at round $t$ is a function of the rewards it has received
  so far and, for a randomised algorithm, of its own random numbers $U_t$:
  \begin{equation*}
    A_t = f_t(R_1, \dots, R_{t-1}, U_t).
  \end{equation*}
  The rewards $R_1, \dots, R_{t-1}$ are entries of the rows $1$ to $t - 1$
  of the table, and the row $t$, $(X_{1, t}, \dots, X_{K, t})$, is
  independent of these rows and of $U_t$. A function of variables that are
  independent of row $t$ is itself independent of row $t$: $A_t$ is
  independent of $(X_{1, t}, \dots, X_{K, t})$. Therefore conditioning on
  $A_t = a$ does not change the distribution of $X_{a, t}$:
  \begin{equation*}
    \E\left[ X_{A_t, t} \mid A_t = a \right] = \E\left[ X_{a, t} \mid A_t =
    a \right] = \E\left[ X_{a, t} \right] = \mu_a,
  \end{equation*}
  hence $\E[X_{\star, t} - X_{A_t, t} \mid A_t] = \mu^\star - \mu_{A_t} =
  \Delta_{A_t} \geq 0$. Independence is essential: an agent that could peek
  at the rewards of round $t$ before choosing, and play the arm with the
  largest one, would get $\E[\max_a X_{a, t}]$ per round, more than
  $\mu^\star$ as soon as another arm can beat the best one at some round,
  and its ``regret'' would be negative. Summing the conditional expectations
  over the rounds gives \cref{prop:regret-decomp}.

  Finally, the comparison is with the best arm \emph{on average}, not with
  the arm whose rewards happened to add up highest. The latter,
  $\E[\max_a \sum_t X_{a, t}] - \E[\sum_t R_t]$, is larger, because a
  maximum over random sums profits from luck; it is used in
  \emph{adversarial} bandits, where the rewards are not drawn from fixed
  distributions.
\end{remark}

\begin{proposition}{Regret decomposition}{regret-decomp}
  $\mathcal{R}_T = \sum_{a=1}^{K} \Delta_a\, \E\left[N_a(T)\right]$.
\end{proposition}

\begin{proof}
  Given $A_t$, the expected reward is $\mu_{A_t}$, so $\E[\sum_t R_t] =
  \E[\sum_t \mu_{A_t}] = \E[\sum_a \mu_a N_a(T)]$. Since $\sum_a N_a(T) =
  T$, $T \mu^\star - \sum_a \mu_a N_a(T) = \sum_a (\mu^\star - \mu_a)
  N_a(T)$; take expectations.
\end{proof}

\begin{intuition}
  The regret is a bill: each play of arm $a$ costs $\Delta_a$ on average,
  and the bill is the number of plays of each arm times its price. Playing
  an almost-optimal arm costs little; playing a bad arm costs a lot.
\end{intuition}

A good algorithm therefore plays each suboptimal arm rarely, and the more
rarely the larger its gap. The estimate of an arm after $n$ plays is the
sample mean $\hat\mu_a$ of its rewards, which can be updated incrementally
(\cref{prop:incremental}).

\section{Greedy and $\varepsilon$-greedy}

The \emph{greedy} algorithm plays each arm once, then always the arm of
largest $\hat\mu_a$. If the best arm happens to give a bad first reward, it
may never be played again: the regret grows linearly with $T$ with positive
probability. The \emph{$\varepsilon$-greedy} algorithm plays a uniformly
random arm with probability $\varepsilon$ and the greedy arm otherwise.

\begin{proposition}{Constant exploration has linear regret}{eps-linear}
  With a constant $\varepsilon > 0$, the regret of $\varepsilon$-greedy
  satisfies $\mathcal{R}_T \geq \frac{\varepsilon\, T}{K} \sum_a
  \Delta_a$.
\end{proposition}

\begin{proof}
  At each round, with probability at least $\varepsilon / K$ arm $a$ is
  played as a random arm, so $\E[N_a(T)] \geq \varepsilon T / K$; apply
  \cref{prop:regret-decomp}.
\end{proof}

Decreasing $\varepsilon_t$ (e.g.\ $\varepsilon_t \propto 1/t$) can give a
logarithmic regret, but the right schedule depends on the unknown gaps.
The next algorithms explore \emph{where it is needed}: arms that are
clearly worse are abandoned quickly, arms that might be the best are
explored more.

\section{Upper confidence bounds}

The principle of \emph{optimism in the face of uncertainty}: play the arm
whose mean \emph{could} be the largest, given the data. It relies on a
concentration inequality.

\begin{lemma}{Hoeffding's inequality (stated without proof)}{hoeffding}
  Let $X_1, \dots, X_n$ be independent random variables with values in
  $[0, 1]$ and mean $\mu$, and $\bar X_n = \frac{1}{n} \sum_i X_i$. For
  every $c > 0$,
  \begin{equation}
    \Prob\left(\bar X_n \geq \mu + c\right) \leq e^{-2 n c^2}, \qquad
    \Prob\left(\bar X_n \leq \mu - c\right) \leq e^{-2 n c^2}.
  \end{equation}
\end{lemma}

\begin{intuition}
  A sample mean rarely strays far from the true mean, and the probability
  of a deviation $c$ falls exponentially in $n c^2$. Turned around: with
  probability at least $1 - \delta$, $\abs{\bar X_n - \mu} \leq
  \sqrt{\ln(2/\delta) / (2n)}$. Deviations of order $1/\sqrt{n}$ are
  typical: four times more data halve the uncertainty. The lemma gives an
  error bar around each estimate.
\end{intuition}

\begin{definition}{UCB1}{ucb1}
  For rewards in $[0, 1]$, UCB1 (Auer, Cesa-Bianchi and Fischer, 2002)
  plays each arm once, then at round $t$ the arm maximising
  \begin{equation}
    \label{eq:ucb1}
    \hat\mu_a + \sqrt{\frac{2 \ln t}{N_a(t-1)}}.
  \end{equation}
\end{definition}

The bonus is large for arms played rarely, and decreases as their estimate
becomes reliable; the slow growth of $\ln t$ ensures that every arm is
eventually retried.

\begin{intuition}
  The index of an arm is the most optimistic mean that is still plausible
  given its data: the estimate plus its error bar. The width of the bar
  comes from Hoeffding's inequality: choosing $c$ so that $e^{-2 n c^2} =
  t^{-4}$ gives $c = \sqrt{2 \ln t / n}$, so the true mean exceeds the
  index only with a tiny probability. The rule corrects itself. An arm with
  a large index is either good, and playing it is right, or poorly known,
  and playing it shrinks its error bar until the arm is seen for what it
  is. Exploration stops by itself where it is no longer needed.
\end{intuition}

\begin{theorem}{Regret of UCB1}{ucb-regret}
  For rewards in $[0, 1]$ and every suboptimal arm $a$ ($\Delta_a > 0$),
  \begin{equation}
    \E\left[N_a(T)\right] \leq \frac{8 \ln T}{\Delta_a^2} + 1 +
    \frac{\pi^2}{3},
    \qquad \text{hence} \qquad
    \mathcal{R}_T \leq \sum_{a : \Delta_a > 0} \left( \frac{8 \ln
    T}{\Delta_a} + \left(1 + \frac{\pi^2}{3}\right) \Delta_a \right).
  \end{equation}
\end{theorem}

\begin{proof}
  Fix a suboptimal arm $a$ and an optimal arm $\star$. Let $c_{t, s} =
  \sqrt{2 \ln t / s}$, let $\hat\mu_{a, s}$ be the mean of the first $s$
  rewards of arm $a$, and $\ell = \lceil 8 \ln T / \Delta_a^2 \rceil$.
  Arm $a$ is played at most $\ell$ times while $N_a < \ell$, so
  \begin{equation*}
    N_a(T) \leq \ell + \sum_{t = K+1}^{T} \ind\left[A_t = a,\
    N_a(t-1) \geq \ell\right],
  \end{equation*}
  where $\ind[E]$ denotes the indicator of an event $E$: $1$ if $E$ occurs,
  $0$ otherwise. Each round that plays $a$ after its first $\ell$ plays adds
  $1$ to the sum. If $A_t = a$, the index of $a$ is at least that of $\star$:
  $\hat\mu_{\star, s} + c_{t, s} \leq \hat\mu_{a, s_a} + c_{t, s_a}$
  for the current counts $s = N_\star(t-1)$ and $s_a = N_a(t-1) \geq
  \ell$. This inequality implies one of the three events
  \begin{equation*}
    \text{(E1) } \hat\mu_{\star, s} \leq \mu^\star - c_{t, s}, \qquad
    \text{(E2) } \hat\mu_{a, s_a} \geq \mu_a + c_{t, s_a}, \qquad
    \text{(E3) } \mu^\star < \mu_a + 2 c_{t, s_a},
  \end{equation*}
  for if none holds, $\hat\mu_{\star, s} + c_{t, s} > \mu^\star \geq
  \mu_a + 2 c_{t, s_a} > \hat\mu_{a, s_a} + c_{t, s_a}$. Event (E3) is
  impossible: $t \leq T$ and $s_a \geq \ell \geq 8 \ln T / \Delta_a^2$
  give $2 c_{t, s_a} \leq 2 \sqrt{2 \ln T \cdot \Delta_a^2 / (8 \ln T)}
  = \Delta_a$. By Hoeffding's inequality, for fixed $s$ and $s_a$, (E1) and
  (E2) have probability at most $e^{-2 s \cdot 2 \ln t / s} = t^{-4}$.
  Since the counts are random, we take a union bound over their possible
  values $1 \leq s, s_a \leq t$:
  \begin{equation*}
    \E[N_a(T)] \leq \ell + \sum_{t=K+1}^{T} \sum_{s=1}^{t}
    \sum_{s_a=1}^{t} 2\, t^{-4}
    \leq \ell + 2 \sum_{t \geq 1} t^{-2} = \ell + \frac{\pi^2}{3},
  \end{equation*}
  and $\ell \leq 8 \ln T / \Delta_a^2 + 1$. The regret bound follows from
  \cref{prop:regret-decomp}.
\end{proof}

\begin{intuition}
  To tell apart two means separated by $\Delta_a$, the error bars must be
  narrower than $\Delta_a$, which takes about $1 / \Delta_a^2$ samples
  since they shrink like $1 / \sqrt{n}$. The factor $\ln T$ is the price of
  confidence: to be wrong with probability about $1/T$ only, the bars must
  be widened by $\sqrt{\ln T}$. Close arms are played often but each play
  costs little; the regret of arm $a$ is $\ln T / \Delta_a$, large only when
  the gap is small, and it never exceeds the trivial bound $\Delta_a T$.
\end{intuition}

The regret of UCB1 grows like $\ln T$, and this is optimal up to the
constant:

\begin{theorem}{Lai--Robbins lower bound (stated without proof)}{lai-robbins}
  For any algorithm whose regret is $o(T^\alpha)$ for every $\alpha > 0$ on
  every bandit of a given parametric family, and every suboptimal arm $a$,
  \begin{equation}
    \liminf_{T \to \infty} \frac{\E[N_a(T)]}{\ln T} \geq
    \frac{1}{\mathrm{KL}(\nu_a \,\|\, \nu_\star)},
  \end{equation}
  where KL is the Kullback--Leibler divergence (Lai and Robbins, 1985).
\end{theorem}

\begin{intuition}
  $\mathrm{KL}(\nu_a \,\|\, \nu_\star)$ measures how much evidence one
  reward of arm $a$ brings that $a$ is not the best arm. Any algorithm that
  does well on every bandit must gather about $\ln T$ units of this
  evidence against each suboptimal arm, hence play it about $\ln T /
  \mathrm{KL}$ times. Arms whose rewards look alike are hard to rule out.
\end{intuition}

\section{Thompson sampling}

Thompson sampling (1933) is Bayesian: maintain a posterior distribution
over the mean of each arm, sample one value from each posterior, and play
the arm whose sample is largest. Each arm is played with the posterior
probability that it is the best.

\begin{intuition}
  An arm that is well known has a narrow posterior, and its sample is close
  to its estimate. An arm that is poorly known has a wide posterior, and
  now and then its sample is large, so it gets tried. An arm that is
  surely worse almost never wins the draw. The randomness of the samples
  plays the role of the error bars of UCB, and it is calibrated by the
  data instead of by a worst-case formula.
\end{intuition}

\begin{algobox}{Thompson sampling for Bernoulli rewards}
  \begin{algorithmic}[1]
    \State $S_a \gets 0$, $F_a \gets 0$ for every arm
      \Comment{successes and failures}
    \For{$t = 1, \dots, T$}
      \State sample $\theta_a \sim \mathrm{Beta}(1 + S_a, 1 + F_a)$ for
        every arm
      \State play $A_t = \argmax_a \theta_a$, observe $R_t \in \{0, 1\}$
      \State $S_{A_t} \gets S_{A_t} + R_t$;\quad
        $F_{A_t} \gets F_{A_t} + 1 - R_t$
    \EndFor
  \end{algorithmic}
\end{algobox}

The $\mathrm{Beta}(1 + S, 1 + F)$ distribution is the posterior of a
Bernoulli mean under a uniform prior after $S$ successes and $F$
failures. Thompson sampling attains the Lai--Robbins bound for Bernoulli
bandits (Kaufmann, Korda and Munos, 2012; Agrawal and Goyal, 2012), and it
is often the best performer in practice.

\begin{remark}{A testbed}{bandit-testbed}
  The script \code{Course/code/bandits.py} compares the algorithms on
  Bernoulli bandits with $K = 10$ arms whose means are drawn uniformly in
  $[0, 1]$, averaging the regret over $100$ such bandits:
  \begin{center}
    \begin{tabular}{@{}lcccc@{}}
      \toprule
      Regret after $T$ rounds & $100$ & $1\,000$ & $10\,000$ &
        $100\,000$ \\
      \midrule
      Greedy & $8.8$ & $28.2$ & $218.6$ & $2121.9$ \\
      $\varepsilon$-greedy, $\varepsilon = 0.1$ & $12.2$ & $57.1$ & $441.0$ &
        $4177.5$ \\
      $\varepsilon$-greedy, $\varepsilon = 0.01$ & $9.8$ & $36.5$ & $227.2$ &
        $857.0$ \\
      UCB1 & $27.6$ & $132.8$ & $362.1$ & $692.9$ \\
      Thompson sampling & $13.9$ & $27.4$ & $44.6$ & $64.1$ \\
      \bottomrule
    \end{tabular}
  \end{center}
  Greedy and $\varepsilon$-greedy have a linear regret: each tenfold
  increase of $T$ multiplies it by about $10$ (by $4$ for $\varepsilon =
  0.01$, which is still in transition). Thompson sampling adds about $19$
  per tenfold increase: the logarithmic growth of the theory. UCB1 is also
  sublinear but has a large constant: its bonus $\sqrt{2 \ln t / n}$,
  designed for the worst case of rewards in $[0, 1]$, keeps exploring
  long after the arms are well estimated, so it is overtaken by
  $\varepsilon = 0.01$ at first and only passes it after $10\,000$ rounds.
  Guarantees of order $\ln T$ say nothing about the constants.
\end{remark}

\section{Gradient bandits}

Another approach learns \emph{preferences} $H(a)$ and plays the softmax
policy $\pi(a) = e^{H(a)} / \sum_b e^{H(b)}$. After playing $A_t$ and
receiving $R_t$, the \emph{gradient bandit} algorithm updates
\begin{equation}
  \label{eq:gradient-bandit}
  H(b) \gets H(b) + \alpha (R_t - \bar R_t) \left( \ind[b = A_t] -
  \pi(b) \right) \qquad \text{for every arm } b,
\end{equation}
where $\bar R_t$ is a \emph{baseline}, for example the average of the past
rewards.

\begin{intuition}
  If the reward is better than usual ($R_t > \bar R_t$), the played arm
  becomes more likely and the others less likely, each in proportion to its
  probability; if it is worse, the opposite. The baseline turns ``was this
  reward good?'' into ``was it better than expected?''. Without it, every
  positive reward would push up whatever arm was played.
\end{intuition}

\begin{proposition}{The gradient bandit is stochastic gradient ascent}%
  {grad-bandit}
  Let $J(H) = \sum_a \pi(a) \mu_a$ be the expected reward. For any
  baseline $b$ that does not depend on $A_t$, the expected update of
  \eqref{eq:gradient-bandit} is $\alpha\, \nabla_H J$.
\end{proposition}

\begin{proof}
  By \cref{prop:softmax-score} (proved in \cref{ch:pg}),
  $\partial \pi(a) / \partial H(b) = \pi(a) (\ind[a = b] - \pi(b))$.
  Hence
  \begin{equation*}
    \frac{\partial J}{\partial H(b)} = \sum_a \mu_a\, \pi(a) \left(
    \ind[a = b] - \pi(b) \right)
    = \E\left[ \mu_{A_t} \left( \ind[A_t = b] - \pi(b) \right)
    \right].
  \end{equation*}
  Given $A_t$, $\E[R_t] = \mu_{A_t}$, and the baseline term has mean zero:
  $\E[\ind[A_t = b] - \pi(b)] = \pi(b) - \pi(b) = 0$.
\end{proof}

This is the one-state case of the policy gradient theorem
(\cref{thm:pg}): the baseline reduces the variance without changing the
expected update, exactly as in actor-critic methods.

\section{Contextual bandits}

In a \emph{contextual bandit}, a context $x_t$ (a user, a page, a
situation) is observed before choosing the arm, and the reward depends on
both. It is an MDP whose next state does not depend on the action:
decisions have no long-term consequences, but the policy must generalise
over contexts. LinUCB (Li et al., 2010) assumes $\E[R \mid x, a] =
\theta_a^\top x$, estimates $\theta_a$ by ridge regression with matrix
$M_a = \lambda I + \sum_{t : A_t = a} x_t x_t^\top$, and plays
\begin{equation}
  \argmax_a \left( \hat\theta_a^\top x_t + \beta \sqrt{x_t^\top M_a^{-1}
  x_t} \right),
\end{equation}
the bonus measuring the uncertainty of the estimate in the direction of the
context. Contextual bandits are the standard model of recommendation and
online advertising.

\begin{intuition}
  $\sqrt{x^\top M_a^{-1} x}$ is, up to the noise level, the standard
  deviation of the predicted reward $\hat\theta_a^\top x$. It is small for
  contexts similar to those already seen with arm $a$ (directions where
  $M_a$ is large), and large for new kinds of contexts. It is the error bar
  of UCB, for a linear model.
\end{intuition}

\section{Summary}

\begin{itemize}[nosep]
  \item The regret is the sum over suboptimal arms of the gap times the
    number of plays; constant exploration gives a linear regret.
  \item UCB1 explores optimistically and has a regret in $O(\ln T)$, the
    optimal order (Lai--Robbins); Thompson sampling samples from the
    posterior, is asymptotically optimal and usually better in practice.
  \item The gradient bandit is a one-state policy gradient method, with a
    baseline.
  \item Exploration in MDPs is harder: an action can be informative
    because of where it leads, not only through its immediate reward
    (\cref{ch:exploration}).
\end{itemize}
